% ── Three-track theorem system ──────────────────────────────
%
%  Track 1 — Mathematical derivations
%  Track 2 — Audit methodology (auditprop, auditrule)
%  Track 3 — Phenomenological observations (phenomenon, interpretation)
%
%  Keeping these separate preserves inferential clarity across
%  400+ pages of mixed formal and interpretive writing.

\theoremstyle{plain}
\newtheorem{theorem}{Theorem}[chapter]
\newtheorem{lemma}[theorem]{Lemma}
\newtheorem{proposition}[theorem]{Proposition}
\newtheorem{corollary}[theorem]{Corollary}

\theoremstyle{definition}
\newtheorem{definition}[theorem]{Definition}
\newtheorem{axiom}[theorem]{Axiom}
\newtheorem{protocol}[theorem]{Protocol}

\theoremstyle{remark}
\newtheorem{remark}[theorem]{Remark}
\newtheorem{example}[theorem]{Example}

% Track 2: Audit propositions and rules
\theoremstyle{definition}
\newtheorem{auditprop}{Audit Proposition}[chapter]
\newtheorem{auditrule}[auditprop]{Audit Rule}

% Track 3: Phenomenological observations
\theoremstyle{remark}
\newtheorem{phenomenon}{Phenomenological Observation}[chapter]
\newtheorem{interpretation}[phenomenon]{Interpretation}
